\section{Implementation and Evaluation}
\label{sec:implementation}

\subsection{Implementation}

Lux Node is implemented in Go, chosen for its concurrency primitives (goroutines, channels) and single-binary deployment. Key implementation details:

\begin{description}
    \item[Networking] libp2p-based peer-to-peer layer with Kademlia DHT for peer discovery and multiplexed connections for concurrent chain operation.
    \item[Consensus Engine] Each chain instantiates its own consensus engine (Avalanche DAG or Snowball linear) as a Go interface, enabling hot-swappable consensus per chain.
    \item[Virtual Machines] A VM interface allows chains to define custom state machines. Built-in VMs: AVM (X-Chain, UTXO), CoreVM (C-Chain, EVM), PlatformVM (P-Chain, staking).
    \item[Database] LevelDB for persistent storage with in-memory caching. Each chain maintains independent database namespaces.
    \item[Bootstrapping] New validators sync state from peers via iterative sampling, converging to the correct state through the same metastable dynamics used for consensus.
\end{description}

\subsection{Testnet Evaluation}

We evaluate on a globally distributed testnet with 2,000 validators across 5 continents.

\subsubsection{Finality Latency}

\begin{table}[h]
\centering
\begin{tabular}{@{}lcc@{}}
\toprule
\textbf{Metric} & \textbf{X-Chain (DAG)} & \textbf{C-Chain (Linear)} \\
\midrule
Median finality & 0.85s & 1.2s \\
P95 finality & 1.4s & 2.0s \\
P99 finality & 2.1s & 2.8s \\
\bottomrule
\end{tabular}
\caption{Finality latency with 2,000 validators, $k=20$, $\alpha=0.8$, $\beta=20$.}
\end{table}

\subsubsection{Throughput}

\begin{table}[h]
\centering
\begin{tabular}{@{}lc@{}}
\toprule
\textbf{Chain} & \textbf{Sustained TPS} \\
\midrule
X-Chain (simple transfers) & 4,500 \\
C-Chain (EVM transactions) & 1,200 \\
P-Chain (staking operations) & 500 \\
\bottomrule
\end{tabular}
\caption{Sustained throughput under normal load.}
\end{table}

\subsubsection{Byzantine Resilience}

We inject Byzantine validators controlling 15\% of stake, executing adaptive strategies:
\begin{itemize}
    \item \textbf{Equivocation}: Sending conflicting responses to different queriers.
    \item \textbf{Liveness attack}: Refusing to respond to delay convergence.
    \item \textbf{Adaptive strategy}: Switching colors to maximize disagreement.
\end{itemize}

Under all attack strategies, correct validators converge within 3 seconds (P99), with zero safety violations across $10^6$ consensus instances.
